\section{Related work}
\label{sec:related}

% STAGE 3: rewritten 8 September 2026.

\paragraph{Degrees of autonomy, and where knowledge enters.} LLM-based AutoML
systems can be ordered by how much of the modelling decision they leave to the
model: hyperparameter optimisers use the language model mainly for the cold
start, pipeline composers search jointly over features, model and
hyperparameters, and adaptive planners revise the plan as validation feedback
arrives. Orthogonally, systems differ in where knowledge of libraries and
methods comes from: the parametric memory of the backbone, retrieved material
such as strategic insights mined from competition
write-ups~\citep{kulibaba2025kompeteai}, or API-level knowledge embedded in
the decision process, as MLZero does with semantic and episodic memory that
suppresses hallucinated and outdated API calls~\citep{fang2025mlzero}.
Section~\ref{sec:axis} splits this second axis into two instructions with
different mechanisms: which library to use, and how to train and validate.

\paragraph{Ablations that disagree.} Multi-agent frameworks for data science
competitions are numerous~\citep{li2024autokaggle,trirat2025automlagent,kulibaba2025kompeteai,fang2025mlzero},
and their reported ablations disagree. On MLE-bench, search structure is the
largest single contributor in one system and has no measurable effect in
another, and memory ranks last in one study and level with search in a
second~\citep{toledo2025aira}. Disagreement of that kind is expected when
ablations are run one factor at a time on different backbones and task sets,
and the factors also interact: advanced search policies gain nothing when
restricted to a weaker operator set~\citep{toledo2025aira}. Our design holds
the backbone, the tasks and the environment fixed and crosses the factors, so
their interaction is measured.

\paragraph{Evaluation against tuned baselines.} Agentic AutoML papers that
compare against classical methods almost always compare against defaults. The
closest work in domain is GRACE-DS~\citep{tsymbalov2026graceds}, which
evaluates LLM-driven AutoML agents on tabular tasks under a structured
environment and reports that structuring the interaction compresses the spread
across backbones. Two differences matter. GRACE-DS normalises scores against
an untuned HistGradientBoosting model, a threshold well below the tuned
XGBoost, LightGBM and CatBoost used here, and its eight backbones are frontier
models of 80B to 1.6T parameters, all larger than the dense 31B model used
here. GRACE-DS varies the backbone under one structure; we vary the structure
under one backbone. Section~\ref{sec:modelaxis} measures the backbone effect under our conditions.
On a scale where 0 is the weakest and 1 the strongest of the eighteen published
methods on each task, swapping the backbone under AutoDS-Tools moves the score
by 0.04; the instruction file that ships with AutoDS-Tools moves it by
$\layerWorthNorm$, and the three architectures with that file removed lie
within $\spanArch$ of each other (Section~\ref{sec:halfrange}).

\paragraph{Scaffolding against model capability.} One study in the embodied setting asked whether scaffolding can substitute for
model capability~\citep{agentspec2026}. It crossed reasoning, memory,
reflection and reinforcement modules with dense models of increasing size and
concluded that external scaffolding can compensate for weaker native
long-horizon reasoning. We follow that design on supervised tabular tasks, where eighteen tuned
methods, among them XGBoost, LightGBM and CatBoost, give a strong reference. Agentic coding has
produced parallel evidence that the harness is a large hidden variable. Across
configurable harnesses evaluated identically, aggregate score spans $23.8$
points~\citep{harnessbench2026}. Harness choice induces up to a 40$\times$
difference in tokens per solved task while the model-to-model pass-rate spread
stays within eight points~\citep{scaffoldeffect2026}. That literature varies
the harness with the task family fixed; we fix the task family and vary how
much of the pipeline the architecture decides in advance, from FEDOT.LLM
through AutoDS-Tools to Terminus-2, and tuned methods on the same split turn
the spread into a fraction of the range between the weakest and the strongest
published method. A third result bears on our design: when a coding harness
is evolved automatically, ablations localise the gain to tools, middleware and
long-term memory, with the system prompt contributing
little~\citep{harness2026prompt}. Our prescribed-knowledge layer is a
prompt-level intervention; Section~\ref{sec:axis} measures its size and
Sections~\ref{sec:grid} and~\ref{sec:kdiscterm} its effect, and the two
instructions in it gain less together than the sum of their separate gains,
as that result predicts.
